% p1-05-null-models

\section{P1 §5. Null Models}

5. Null Models

    Figure (fig-p1-ch5-null-models). Null-models triptych --- log-uniform / permutation / randomized
                                          exponent grammar.

  5.1 Rationale for Explicit Null Models
  Without a properly specified null distribution, observed residuals dC(Xi) are statistically
  uninterpretable. The grammar G$\varphi$ is sufficiently expressive that small residuals are
  guaranteed to exist for large C, regardless of whether the target values carry algebraic
  structure. The null models defined here answer: how small would residuals be if the
  target values had no structural relation to the grammar?

  Three complementary null models are specified.

  5.2 Null Model M0: Log-Uniform Random Baseline
  Under M0, the target values {Xi}i=1N are replaced by independent samples from a
  log-uniform distribution on the empirical range of T:

  log Xi(0) $\sim$ Uniform[log Xmin, log Xmax], 9

  where [Xmin, Xmax] matches the empirical range. The coverage functional under M0,

  rho0(C, $\varepsilon$) := EX(0) $\sim$ M0[(1)/(N)sumi=1N 1[dC(Xi(0)) < $\varepsilon$]], 10

  is computed by Monte Carlo sampling over B $\geq$ 10{,}000 independent realizations of
  T(0). The null hypothesis is rho(C, $\varepsilon$) = rho0(C, $\varepsilon$); the alternative is rho(C, $\varepsilon$) > rho0(C,
  $\varepsilon$).

  5.3 Null Model Mperm: Permutation-Based Baseline
  The permutation null model Mperm preserves the empirical distribution of T while
  destroying any structural relation between individual constants and the grammar.

  Under Mperm, the target set is replaced by a random permutation of its elements
  combined with a random sign-flip in log-space:

  Tpi := {exp(± log Xpi(i))}i=1N, pi uniform on SN. 11

  The coverage functional under Mperm is

  rhoperm(C, $\varepsilon$) := Epi $\sim$ SN[(1)/(N)sumi=1N 1[dC(Xpi(i)) < $\varepsilon$]]. 12

  Because Mperm samples from the same empirical distribution as T, it controls for
  distributional effects without assuming log-uniformity, making it more conservative
  than M0.

  5.4 Null Model Mrand: Randomized Exponent Grammar
  A third control uses a randomized grammar Mrand in which exponent vectors are
  drawn uniformly from the same bounded integer domain as S(C), but with basis vectors
  replaced by independent random positive reals:

  k, p, m, q $\sim$ Unif({-C,…,C}4), n $\sim$ Unif({1,…,9}), log vrand $\sim$ Unif[0,1]4. 13

  This generates an ensemble Srand(C) with the same cardinality as S(C) but without the
  specific algebraic structure of the $\varphi$-basis. Comparing coverage under G$\varphi$ to coverage
  under Mrand isolates the contribution of the specific algebraic structure.

  5.5 Summary of Null Model Hierarchy
  The three null models form a hierarchy of increasing conservatism:

  M0 $\subset$ Mperm $\subset$ Mrand.

  A compressibility claim has statistical content only if it survives all three comparisons
  simultaneously.

  5.6 Negative Controls: Randomly Rotated Constants
  In addition to the three null models, a battery of negative controls is constructed using
  randomly rotated versions of T:

  Tneg(j) = {Xi $\cdot$ Uj}i=1N, Uj $\sim$ Log-Uniform(0.5, 2.0), 13.a

  for j = 1, …, 50 independent realizations. Coverage under G$\varphi$ applied to these rotated
  datasets should not exceed the null expectation, because the rotation destroys any
  algebraic alignment. Failure of this check would indicate that the grammar is simply
  dense in the relevant range, not genuinely selective.

